\documentclass[conference]{IEEEtran}
\usepackage{cite}
\usepackage{amsmath,amssymb,amsfonts}
\usepackage{graphicx}
\usepackage{textcomp}
\usepackage{xcolor}
\usepackage{hyperref}

\title{Synthetic Environment Impact on Object Detection Generalization\\
in Robotic Perception Systems}

\author{\IEEEauthorblockN{Sahan}
\IEEEauthorblockA{Department of Computer Engineering\\
University of Ruhuna, Sri Lanka}}

\begin{document}
\maketitle

\begin{abstract}
% TODO Week 1: Fill in after setup complete
This paper investigates how the structure of synthetic training environments
affects the generalization of object detection models in robotic perception tasks.
Using ROS2 Jazzy, Gazebo Harmonic, and YOLOv8, we construct three distinct
simulation environments and train separate models on each.
A 9-combination cross-evaluation benchmark reveals domain shift patterns.
We further demonstrate few-shot domain adaptation by fine-tuning across environments.
A weighted scoring formula provides deployment recommendations.
Findings suggest that [FILL AFTER RESULTS].
\end{abstract}

\section{Introduction}
% TODO Week 1

\section{Related Work}
% TODO Week 2

\section{Methodology}
% TODO Weeks 3-5

\section{Results}
% TODO Week 6

\section{Discussion}
% TODO Week 7

\section{Conclusion}
% TODO Week 7

\bibliography{references}
\bibliographystyle{IEEEtran}

\end{document}
